\section{Architettura del sistema}

\subsection{Diagramma a blocchi}
\begin{figure}[H]
    \centering
    % \includegraphics[width=\textwidth]{architettura-blocchi.png}
    \fbox{\parbox{0.85\textwidth}{\centering\vspace{1em}
        \TODO{Inserire diagramma a blocchi dell'architettura. Componenti:
        Mobile App (Expo/RN) $\leftrightarrow$ HTTPS/JSON $\leftrightarrow$
        Flask Backend (Routes $\to$ Services $\to$ Models $\to$ SQLAlchemy)
        $\leftrightarrow$ SQLite.}
        \vspace{1em}}}
    \caption{Architettura a blocchi del sistema.}
    \label{fig:arch}
\end{figure}

\subsection{Suddivisione in strati del backend}
Il backend è organizzato secondo il pattern \emph{Routes $\to$ Services $\to$
Models}, con utility trasversali.

\begin{table}[H]
    \centering
    \begin{tabularx}{\textwidth}{@{}lX@{}}
        \toprule
        \textbf{Strato} & \textbf{Responsabilità} \\
        \midrule
        \code{app/routes/} &
        Punti di ingresso HTTP. Parsing del payload, dispatch ai servizi,
        formattazione della risposta JSON, mapping degli errori in codici
        di stato. \\
        \code{app/services/} &
        Logica di business. Transazioni, eccezioni di dominio
        (\code{PublishError}, \code{AuthError}), verifica PoW, controlli di
        autorizzazione. \\
        \code{app/models/} &
        Definizione delle classi SQLAlchemy: \code{Artist}, \code{Album},
        \code{Song}, \code{User}, \code{PowChallenge}. Metodi di
        serializzazione (\code{to\_dict}, \code{to\_lrclib}). \\
        \code{app/utils/} &
        Helper trasversali: encoder/decoder LRC, utility crittografiche
        (token, hashing). \\
        \code{app/config.py} &
        Configurazione (env-aware): URI database, chiavi JWT, scadenze
        token, limiti dimensionali, difficoltà PoW. \\
        \code{app/extensions.py} &
        Singleton delle estensioni Flask (db, jwt, cors, migrate). \\
        \code{main.py} &
        Entrypoint: app factory, seeding del catalogo iniziale, avvio del
        server di sviluppo. \\
        \bottomrule
    \end{tabularx}
    \caption{Layout del backend.}
\end{table}

\subsection{Suddivisione del frontend mobile}
\begin{table}[H]
    \centering
    \begin{tabularx}{\textwidth}{@{}lX@{}}
        \toprule
        \textbf{Cartella/File} & \textbf{Ruolo} \\
        \midrule
        \code{app/} & Schermate (file-based routing di expo-router). \\
        \code{app/\_layout.tsx} & Root layout, registra \code{AuthProvider}. \\
        \code{app/index.tsx} & Home con stato di connessione e pillola profilo. \\
        \code{app/search.tsx} & Ricerca interattiva. \\
        \code{app/explore.tsx} & Catalogo paginato con ordinamento. \\
        \code{app/song/[id].tsx} & Dettaglio brano con testo piano e sincronizzato. \\
        \code{app/publish.tsx} & Form di pubblicazione (JWT o PoW). \\
        \code{app/login.tsx}, \code{register.tsx} & Schermate auth. \\
        \code{app/my-lyrics.tsx} & Elenco delle pubblicazioni dell'utente. \\
        \code{app/edit-song/[id].tsx} & Form di modifica. \\
        \code{app/admin.tsx} & Pannello amministratore (solo utenti con \code{is\_admin}). \\
        \code{src/api.ts} & Client REST per tutte le chiamate al backend (auth, lyrics, admin). \\
        \code{src/AuthContext.tsx} & Stato globale dell'autenticazione. \\
        \code{src/storage.ts} & Wrapper cross-platform per la persistenza dei token. \\
        \code{src/dialog.ts} & Alert/confirm cross-platform (mobile e web). \\
        \code{src/sha256.ts} & Implementazione client di SHA-256 per la PoW. \\
        \bottomrule
    \end{tabularx}
    \caption{Layout del frontend mobile.}
\end{table}

\subsection{Schema del database}
\begin{figure}[H]
    \centering
    % \includegraphics[width=\textwidth]{schema-db.png}
    \fbox{\parbox{0.85\textwidth}{\centering\vspace{1em}
        \TODO{Inserire diagramma ER del database (entità: users, artists,
        albums, songs, pow\_challenges; chiavi e foreign key).}
        \vspace{1em}}}
    \caption{Schema entità-relazione del database.}
    \label{fig:schema-db}
\end{figure}

\subsubsection{Tabelle principali}
\begin{itemize}
    \item \code{users} \emph{(id, username, email, password\_hash, is\_admin, created\_at)} ---
    utenti registrati. Username ed email univoci (case-insensitive).
    \code{is\_admin} (booleano, default \code{false}) abilita le funzioni
    di moderazione.
    \item \code{artists} \emph{(id, name)} --- artisti, nome univoco.
    \item \code{albums} \emph{(id, title, artist\_id)} ---
    album opzionalmente associati a un brano.
    \item \code{songs} \emph{(id, title, artist\_id, album\_id, lyrics,
    synced\_lyrics, duration, instrumental, user\_id, created\_at)} ---
    brani con testo piano e sincronizzato. \code{user\_id} nullable: \code{NULL}
    indica un brano anonimo (pubblicato via PoW).
    \item \code{pow\_challenges} \emph{(id, token, difficulty, used, created\_at)} ---
    sfide PoW emesse e consumate.
\end{itemize}
